% Standalone wrapper that builds the methods and results sections as one PDF.
%   latexmk -pdf compile_sections.tex
% Inputs: methods_L1_L2.tex, results_validation.tex. Figures are read from
% ../outputs/validation/Ruby2D/, which is not in the repository (outputs/ is
% gitignored), so a fresh clone cannot rebuild this PDF; the committed PDF is
% the reference copy.
%
% Status (2026-10-15): the Methods section was checked against the code for
% catalog version 2026-10-14.1. The Results section is a historical record of
% analyses run April to July 2026 and was not recomputed. Audit:
% docs/doc_audit_2026-10-15/audit_compile_sections.md.

\documentclass[11pt]{article}
\usepackage[margin=1in]{geometry}
\usepackage{amsmath}
\usepackage{booktabs}
\usepackage{graphicx}
\graphicspath{{../outputs/validation/}}
\usepackage[utf8]{inputenc}
\usepackage[T1]{fontenc}
\usepackage{siunitx}
\sisetup{range-phrase = --, range-units = single}
\DeclareSIUnit{\dbar}{dbar}
\usepackage{hyperref}
\setlength{\emergencystretch}{3em}
\sloppy
\hypersetup{colorlinks=true, linkcolor=blue, urlcolor=blue}

% Visible status note: a boxed paragraph that marks a change, a correction,
% or a historical/retracted passage, so the marker survives into the PDF.
\newcommand{\statusnote}[1]{%
  \par\smallskip\noindent
  \fbox{\parbox{\dimexpr\linewidth-2\fboxsep-2\fboxrule\relax}{\small #1}}%
  \par\smallskip}

\title{PUV Pipeline: Methods and Validation Results}
\author{Holden Leslie-Bole \\ Scripps Institution of Oceanography}
\date{Draft of April 2026; methods updated 15 October 2026}

\begin{document}
\maketitle

\section*{How to read this document}

This is a working draft of the methods and results sections written for the
first PUV paper. It is kept because it is the most complete written account of
the L1 to L3 processing and of the validation tests run in spring and summer
2026.

\begin{itemize}
  \item \textbf{Methods} (Section~\ref{sec:methods}) were checked against the code
    on 2026-10-15 and describe the processing that produced catalog version
    2026-10-14.1. Each paragraph that changed carries a boxed note saying what
    changed and when.
  \item \textbf{Results} (Section~\ref{sec:results}) record analyses run between
    April and July 2026 on the catalog as it then stood (33 to 38 instruments,
    17-minute segments for the earliest). None was recomputed for this revision.
    Passages later retracted are boxed and labeled.
\end{itemize}

For the current pipeline, start with \texttt{README.md}, then
\texttt{docs/pipeline\_levels.md} (status by level),
\texttt{PIPELINE\_NOTES.md} (design reasons) and \texttt{CHANGELOG.md}. The
long-form manual is \texttt{docs/PUV\_Pipeline\_Guide.pdf}. The spectral-estimator
choice is documented in \texttt{docs/multitaper\_writeup.pdf}, and the faults of
the legacy processing code in \texttt{docs/legacy\_pipeline\_defects.pdf}. Current
model-comparison findings live in the separate model-comparison paper repository (not public), not here.

\tableofcontents
\newpage

\section{Methods}
\label{sec:methods}
% PUV Pipeline: methods section, \input by compile_sections.tex.
% Methods text checked against the code on 2026-10-15 (catalog version
% 2026-10-14.1). Paragraphs changed in that pass carry a visible
% "Changed 2026-10" note. The audit behind each change is in
% docs/doc_audit_2026-10-15/audit_compile_sections.md.
% Requires the \statusnote macro defined in compile_sections.tex.

\subsection{Nearshore wave measurements}

Waves and currents were measured with bottom-mounted Nortek Vector acoustic
Doppler velocimeters (ADVs) recording pressure and velocity (PUV) continuously at
\SI{2}{\hertz}. The catalog holds 65 instrument-records from 46 deployments on the
San Diego County coast and at Catalina Island, 2014--2026, at nominal depths of
\SIrange{5}{30}{\meter} (median \SI{7.5}{\meter} over the 72 configured records).
Deployment metadata (instrument heading, clock drift, surveyed position, and the
pressure-port height above the bed, $z_s$, called \texttt{doffp} in the code)
come from the field deployment notes. Across the 72 configured records $z_s$
ranges from \SIrange{0.38}{1.09}{\meter}.

\statusnote{Changed 2026-10. The April draft described the TBR23 array only
(nominal depths \SIrange{5}{7}{\meter}, $z_s = $ \SIrange{0.75}{0.77}{\meter}).
The ranges above are read from the deployment configs through
\texttt{config/deployment\_registry.m}.}

\subsection{Level 1: Quality control and coordinate rotation}

Raw data are read either from the Nortek ExploreV ASCII export
(\texttt{.dat}/\texttt{.sen}/\texttt{.hdr}) or from the recorder binary
(\texttt{.VEC}) by \texttt{L1\_raw\_to\_qc/read\_VEC.m}. Each recorder file is
treated as one burst. Clock drift measured at recovery is removed linearly over
the deployment. For the pre-2019 recorders, whose internal clock ran from a wrong
epoch, time is recovered from the \texttt{MMDDHHMM} file names; that clock was
Pacific local time, so a fixed \texttt{instr.clockOffsetHours} is added to
reach UTC (7 h for the 2014 records, 8 h for the later ones; each checked against
the tide with \texttt{scripts/audit\_clock\_lag.m}). A gap between consecutive sensor records longer than
\texttt{cutoffGapSec} (\SI{1}{\second} by default, \SI{60}{\second} for the
pre-2019 recorders) marks battery depletion, and all data after the first such
gap are discarded. The gap test runs within each burst, so a seam between files
is not read as a cutoff.

Quality control is applied per channel (\texttt{L1\_raw\_to\_qc/puv\_channel\_qc.m},
thresholds from \texttt{l1\_qc\_opts.m}), so a failed thermistor or pressure
sensor does not discard good velocity:
\begin{itemize}
  \item \emph{Velocity} fails where the minimum beam correlation is below 70\%,
    where the frame is toppled ($|\text{pitch}|$ or $|\text{roll}| \ge
    \SI{30}{\degree}$), or where the tilt is unstable: the standard deviation of
    pitch or roll over a \SI{60}{\second} (120-sample) moving window exceeds
    \SI{2}{\degree}.
  \item \emph{Pressure} fails outside $[p_\text{ref}/2,\; 2p_\text{ref}]$, where
    $p_\text{ref}$ is the median of the in-water samples ($p > \SI{0.5}{dbar}$) in
    the first 10\% of the record (at least 1000 samples). A reference taken from
    the start of the record survives a sensor that saturates later in the
    deployment, which a whole-record median does not.
  \item \emph{Temperature} fails outside the site's plausible range
    \texttt{cfg.qcOpts.Tvalid}, set to \SIrange{9}{26}{\celsius} in every catalog
    config (the code default, \SIrange{-2}{40}{\celsius}, only catches gross
    failure). Where the thermistor is valid and the instrument's sound speed
    differs from the one implied by the measured temperature by more than 0.2\%,
    velocity is rescaled and the segment is marked suspect downstream.
\end{itemize}
Failed samples are set to NaN in their own channel only.

\statusnote{Changed 2026-10. The April draft described a fixed
$\pm\SI{5}{\degree}$ pitch and roll limit, a whole-deployment median pressure
reference, and row-level NaN-ing of every channel. The per-channel QC and the
first-window pressure reference replaced those in July 2026
(\texttt{PIPELINE\_NOTES.md}, \texttt{docs/OUTSTANDING\_channel\_decoupling.md}).
A vertical-transit (off-bed) pressure test exists in \texttt{puv\_channel\_qc.m}
but is off in every config.}

Samples with a stable but non-zero tilt (for example from a pipe bent during a
storm) are kept and rotated to the local vertical sample by sample, using the
measured pitch $\phi$ and roll $\theta$, whenever the median tilt exceeds
\SI{1}{\degree}. The code applies roll first and then pitch,
%
\begin{equation}
  \mathbf{v}_\text{corrected} =
  \mathbf{R}_\text{pitch}(\phi)\;\mathbf{R}_\text{roll}(\theta)\;
  \mathbf{v}_\text{raw},
  \label{eq:tilt}
\end{equation}
%
where
\[
  \mathbf{R}_\text{roll} = \begin{bmatrix}
    1 & 0 & 0 \\
    0 & \cos\theta & -\sin\theta \\
    0 & \sin\theta &  \cos\theta
  \end{bmatrix}, \qquad
  \mathbf{R}_\text{pitch} = \begin{bmatrix}
    \cos\phi & 0 & \sin\phi \\
    0 & 1 & 0 \\
    -\sin\phi & 0 & \cos\phi
  \end{bmatrix}.
\]

\statusnote{Corrected 2026-10. Earlier versions of this section, and the
comment above the loop in \texttt{PUV\_raw\_process.m}, write the product as
$\mathbf{R}_\text{roll}\mathbf{R}_\text{pitch}$. The expanded components in that
loop are those of $\mathbf{R}_\text{pitch}\mathbf{R}_\text{roll}$, which
\eqref{eq:tilt} now states. Whether that order matches Nortek's pitch and roll
convention has not been checked here.}

The tilt-corrected velocity is rotated to a geographic buoy frame ($+x$ west,
$+y$ north, $+z$ up, the left-handed frame used by CDIP MOP). For records in
instrument (XYZ) coordinates the rotation uses the heading, either the median of
the onboard compass or, where the compass drifts or steps at constant tilt, the
mount heading logged in the field, plus the magnetic declination from the
International Geomagnetic Reference Field (\texttt{igrfmagm}, IGRF-14 where the
installed release supports it, otherwise IGRF-13 capped at 2025.0) at the
deployment position and start date. Records stored in Nortek ENU coordinates,
which are magnetic, have the same declination added
(\texttt{enu\_to\_buoy\_frame.m}).

\statusnote{Changed 2026-10. The declination for ENU records was added on
2026-10-07 (about \SI{11.07}{\degree} at San Diego; SOL24, LPL24, LPL25A
rebuilt). The rule for preferring the logged heading dates from the same day.
Twelve records whose logged and compass headings differ by more than
\SI{10}{\degree} keep the compass value and are listed in
\texttt{shared/heading\_disputed\_records.m}; their directional products carry
at least that uncertainty.}

\subsection{Level 2: Spectral analysis}

\subsubsection{Segmentation and preprocessing}

The L1 time series is divided into non-overlapping \SI{1}{\hour} segments of 7200
samples, starting at the first UTC top of the hour, which matches the hourly
cadence of the CDIP MOP hindcast. Random Doppler error in a segment-mean velocity
falls as $1/\sqrt{N}$, so the 7200-sample mean has a noise floor
$\sqrt{7200/2048} \approx 1.87$ times lower than the 2048-sample (17-minute)
segments used before May 2026. Velocity and pressure products are validated
separately: a segment's velocity products are kept when no more than 10\% of its
velocity samples are NaN (\texttt{segValid\_vel}), and likewise for pressure
(\texttt{segValid\_p}). Remaining gaps are filled by linear interpolation in the
interior and by holding the nearest valid value at a segment edge. Velocities
are then rotated to a shore-normal frame ($+x$ onshore, $+y$ alongshore north)
using either a shore-normal angle set in the config (\texttt{instr.shorenormal})
or the angle CDIP publishes for the MOP station (\texttt{instr.mopStation}).

\statusnote{Changed 2026-10. Before 2026-10-09 an edge gap was filled by
extrapolating the interior trend, which inflated infragravity and velocity
statistics on segments with edge gaps (Hs moved by at most \SI{3.5}{\centi\meter}
except one segment, $U_b$ and $u_\text{rms}$ by about 4\% median and 25\% at
the 95th percentile on edge-gap segments; \texttt{docs/catalog\_changes/2026-10-13\_combined\_rebuild.md}).
Pressure-only and velocity-only validity replaced a single all-channel test in
July 2026. The April draft also cited a within-hour stationarity audit; that
audit is not reproduced in this repository.}

For each valid segment the mean pressure gives the total depth
\begin{equation}
  H = \frac{\overline{P} \times 10^4}{\rho g} + z_s ,
  \label{eq:depth}
\end{equation}
with $\overline{P}$ the segment-mean pressure in dbar,
$\rho = \SI{1025}{\kilogram\per\meter\cubed}$ and
$g = \SI{9.81}{\meter\per\second\squared}$. Segment means of $u$, $v$, $w$ and
temperature are stored before detrending. Pressure is converted to meters of
water and all channels are linearly detrended.

\subsubsection{Spectral estimation}

Spectra are estimated with the multitaper method of Thomson (1982)
(\texttt{shared/psd\_multitaper.m}) on the whole 7200-sample segment, with
time--bandwidth product $NW = 4$. The code keeps the Slepian tapers whose
concentration eigenvalue exceeds 0.99, which for $NW = 4$ is $K = 6$ (the
seventh has eigenvalue 0.937), giving about $2K = 12$ degrees of freedom per
frequency bin. Bins are spaced $\Delta f = f_s/N = \SI{2.78e-4}{\hertz}$; the
resolution bandwidth is $2NW/T = \SI{2.2e-3}{\hertz}$ for $T = \SI{3600}{\second}$,
so neighboring bins are not independent. Auto- and cross-spectra among $p$, $u$
and $v$ use the same tapers, which keeps the directional moments consistent.
For comparison, the Welch estimator used before April 2026 (Hann window,
\texttt{nfft} = 256) had bins spaced \SI{7.8e-3}{\hertz}.

\statusnote{Corrected 2026-10. Earlier versions of this document, the multitaper
note and several repository docs state $K = 7$ tapers and 14 degrees of freedom.
\texttt{psd\_multitaper.m} has used the eigenvalue $> 0.99$ rule since it was
written (April 2026), which for $NW = 4$ keeps six tapers (eigenvalues from
\texttt{dpss(7200,4)}: 1.000, 1.000, 1.000, 1.000, 0.999, 0.993, 0.937, 0.699).
The April draft also described the resolution of the multitaper estimate as
comparable to Welch with 256-sample sub-segments; at a 1-hour segment it is
finer by a factor of about 3.5 in bin spacing.}

\subsubsection{Pressure correction}

The pressure spectrum is converted to a surface-elevation spectrum with the
linear transfer function
%
\begin{equation}
  K_p(f) = \frac{\cosh(k z_s)}{\cosh(k H)} ,
  \qquad S_\eta(f) = \frac{S_p(f)}{K_p^2(f)} ,
  \label{eq:kp}
\end{equation}
%
where $k(f)$ solves $\omega^2 = g k \tanh(kH)$ by Newton iteration
(\texttt{shared/get\_wavenumber.m}). Bins with $K_p < 0.1$ are set to zero to
avoid amplifying noise at high frequency, and the resulting cutoff frequency
$f_\text{cut}$ is stored per segment.

\subsubsection{Bulk wave parameters}

Significant wave height $H_s = 4\sqrt{m_0}$, with $m_0 = \int S_\eta\,df$, is
reported over the total band (\SIrange{0.004}{0.25}{\hertz}), the sea--swell band
(\SIrange{0.04}{0.25}{\hertz}) and the infragravity band
(\SIrange{0.004}{0.04}{\hertz}). The peak period $T_p = 1/f_p$ is the frequency of
maximum $S_\eta$ within the sea--swell band (argmax, not a centroid). The mean
period $T_{m02} = \sqrt{m_0/m_2}$ and the energy flux
$F = \rho g \int c_g S_\eta\,df$, with $c_g$ the linear group velocity, are both
evaluated over the total band. The mean wave direction is
$\operatorname{atan2}(\bar b_1, \bar a_1)$, where $\bar a_1$ and $\bar b_1$ are the
first-harmonic directional moments from the pressure--velocity cross-spectra,
weighted by $S_\eta$ over the sea--swell band (\texttt{shared/compute\_bulk\_params.m}).

\subsubsection{Near-bed velocity and stress}

Near-bed orbital velocity is obtained by scaling each frequency bin of the
velocity FFT from the sensor height to the bed. Under linear theory
$u(z) \propto \cosh(kz)/\sinh(kH)$ with $z$ measured up from the bed, so the
bed-to-sensor ratio is
\begin{equation}
  \frac{u_\text{bed}}{u_\text{sensor}} = \frac{1}{\cosh(k z_s)} ,
  \label{eq:bedvel}
\end{equation}
and the scaled spectrum is inverse transformed (\texttt{shared/bed\_velocity\_ifft.m}).
$U_b$ is the root-mean-square of $\sqrt{u_\text{bed}^2 + v_\text{bed}^2}$ over the
segment, and $A_w$ is the corresponding orbital excursion.

\statusnote{Corrected 2026-10. Earlier versions of this section gave the factor as
$\cosh(kH)/\cosh[k(H - z_s)]$, which is the form the code's own header records as
the original, wrong one. Because $k z_s \approx 0.04$ for swell at
\SI{9}{\meter}, the factor in \eqref{eq:bedvel} differs from 1 by about 0.1\%
there.}

Wave bed shear stress uses the Swart (1974) friction factor with $T_p$ as the
wave period. The Nikuradse roughness is $k_s = 2.5 D_{84}$ where a measured
grain-size distribution exists for the site (\texttt{shared/site\_grain\_size.m});
otherwise the legacy form $k_s = 10 D_{50}$ is used with the placeholder
$D_{50} = \SI{0.25}{\milli\meter}$. As of catalog 2026-10-14.1, 19 of 65 records
use a measured distribution, 4 an assumed site default of \SI{217.4}{\micro\meter}
and 42 the placeholder. Reynolds stresses ($\overline{u'w'}$, $\overline{v'w'}$,
$\overline{u'v'}$) and turbulent kinetic energy come from the detrended
velocities. Velocity skewness $S_k$ and asymmetry $A_s$ (Hilbert form, negative
for the forward-pitched shoaling wave shape) follow Elgar and Guza (1985) and
Hoefel and Elgar (2003).

\statusnote{Corrected 2026-10. The April draft gave $k_s = 2.5 D_{50}$ with a
single \SI{0.25}{\milli\meter} grain size. The grain-size tally is from the L3
provenance sidecars, as quoted in \texttt{README.md}.}

\subsection{Validation against the MOP wave model}

\statusnote{Historical method (April 2026). This subsection describes the
comparison behind the results in Section~\ref{sec:results-status}. The scripts
that implemented it are in \texttt{validation/archive/} (\texttt{compare\_PUV\_MOP.m},
\texttt{run\_PUV\_MOP\_validation.m}). The fourth test below, peak broadening
measured by $Q_p$ on an interpolated grid, produced a result that was retracted on
2026-07-24 as a resolution artifact. The current comparison reduces both spectra
to the model's frequency bins first (\texttt{validation/compare\_shape\_matched.m},
\texttt{shared/shoal\_mop\_to\_site.m}).}

PUV-derived wave parameters were compared with hourly output from the CDIP
Monitoring and Prediction (MOP) system, a spectral wave model initialized from
offshore CDIP buoys that provides $S_\eta(f)$ at the \SI{10}{\meter} contour of
each MOP transect. MOP spectra were shoaled to the PUV depth by conservation of
energy flux,
%
\begin{equation}
  S_\eta(f, h_\text{PUV}) = S_\eta(f, \SI{10}{\meter}) \times
  \frac{c_g(f, \SI{10}{\meter})}{c_g(f, h_\text{PUV})} .
\end{equation}
%
Shoaled MOP wave heights were integrated over the sea--swell band with the
non-uniform MOP bandwidth vector. In April 2026 the PUV segments were
\SI{17}{\minute} long and were interpolated to the MOP hourly grid. Agreement was
summarized by $R^2$, RMSE, bias (PUV $-$ MOP) and the Willmott (1981) skill
score.

The frequency-resolved ratio
$S_\eta^\text{PUV}(f) / S_\eta^\text{MOP,shoaled}(f)$ was averaged over matched
hours in three bands: low swell (\SIrange{0.04}{0.08}{\hertz}), mid swell
(\SIrange{0.08}{0.14}{\hertz}) and high swell (\SIrange{0.14}{0.25}{\hertz}).

Four candidate mechanisms for the systematic difference were tested:
(1)~nonlinear shoaling, through the correlation of the spectral ratio with the
Ursell number $\mathit{Ur} = H_s L_p^2 / h^3$; (2)~directional narrowing during
shoaling, by comparing the first-harmonic spread
$\sigma_1 = \sqrt{2(1 - \sqrt{a_1^2 + b_1^2})}$ between PUV and MOP; (3)~bound
long waves, through the scaling of low-frequency excess with $H_s^2$; and
(4)~spectral peak broadening in the model, through the Goda peakedness
$Q_p = (2/m_0^2) \int f\, S_\eta^2\, df$ and the half-power bandwidth.

\subsection{Level 3: Wave forcing characterization}

\subsubsection{Frequency-band energy decomposition}

The L2 surface-elevation spectrum is integrated over three bands, infragravity
(\SIrange{0.004}{0.04}{\hertz}), swell (\SIrange{0.04}{0.10}{\hertz}) and sea
(\SIrange{0.10}{0.25}{\hertz}), giving a band wave height, energy flux and mean
period for each 1-hour segment, and a flag naming the band with the largest
energy flux.

\subsubsection{Storm event detection}

Storm events are taken from the $H_s$ series as exceedances of
\SI{1.5}{\meter} lasting at least \SI{6}{\hour}, with events closer than
\SI{12}{\hour} merged. These thresholds are fixed in
\texttt{L3\_forcing/PUV\_L3\_storms.m}, not read from the config. To describe the
forcing during PUV data gaps (for example when an instrument was damaged in a
storm), hourly MOP output is loaded for a window extending \SI{7}{days} beyond
each end of the record, and storms found in the MOP record but absent from the
PUV record are flagged as gap-period storms.

\subsubsection{Transport proxies}

Bottom energy flux $F_b = \rho g \int c_g S_\eta\,df$ is integrated over the
sea--swell band and accumulated over the deployment for comparison with surveyed
bed change. The Shields parameter
$\theta = \tau_b / [(\rho_s - \rho) g D_{50}]$ uses the wave stress amplitude
$\tau_b$ from L2, with the critical value from Soulsby (1997)
(Soulsby--Whitehouse), and a segment is flagged as mobilized when
$\theta$ exceeds it. Kinematic viscosity is computed per segment from the
segment's mean temperature (\texttt{shared/seawater\_nu.m}), and the settling
velocity $w_s$ follows Ferguson and Church (2004) with nominal-diameter constants
$C_1 = 20$, $C_2 = 1.1$. The Rouse number $P = w_s / (\kappa u_*)$ uses the
friction velocity of the mean combined wave--current stress $\tau_m$
(Soulsby, 1997), and is also given at $D_{16}$, $D_{50}$ and $D_{84}$ where a
measured distribution exists.

\statusnote{Changed 2026-08 and 2026-10. Until 2026-08-13 the Rouse number used
the oscillatory wave stress amplitude, which overstated suspension (median $P$
4.8 against a corrected 41.5 at TBR23 MOP580\_5m, \texttt{docs/pipeline\_levels.md}).
Settling velocity used smooth-sphere constants ($C_2 = 0.4$) until 2026-08-13 and
$C_2 = 1.0$ until 2026-10-09; viscosity was a fixed \SI{1e-6}{\meter\squared\per\second}
until 2026-10-09. Six settling-velocity formulas could not be distinguished
against the TBR23 surveys; they spread about $\pm$15\% at \SI{250}{\micro\meter}.}

\subsubsection{Tidal current decomposition}

Tidal harmonic analysis is run on the segment-mean velocities with
\texttt{t\_tide} (Pawlowicz et al., 2002), using complex input $u + iv$ so the
cross-shore and alongshore constituents are fit together, and only over the span
where velocity exists. Each call is seeded (\texttt{rng(0,'twister')}) because the
bootstrap error estimate decides which constituents enter the prediction. The
tidal depth reference is the NOAA prediction for La Jolla, Scripps Pier (station
9410230), fetched through the lab toolbox function \texttt{getztide2}; in the
catalog build every record has NOAA depth. Subtidal residual currents (observed
minus tidal) carry the wave-driven mean flows, including undertow and longshore
currents.

For TBR23 the NOAA prediction matched the PUV depth record better than a
\texttt{t\_tide} fit to the PUV pressure ($R = 0.995$ against $0.987$), and the
lag of maximum correlation was 0~h, which confirmed that the TBR23 instrument
clocks were in UTC. The same test later showed that the file-name clocks of the
pre-2019 recorders were Pacific local time (see Level 1);
\texttt{scripts/audit\_clock\_lag.m} repeats it for those records.

\statusnote{Changed 2026-10. The seeding and the velocity-span restriction date
from 2026-10-09. A change set in progress on 2026-10-15 makes the tide station a
per-deployment setting and makes a missing \texttt{t\_tide} fatal; until it
lands, every site uses station 9410230, which is the wrong tide for a site far
from La Jolla. See \texttt{CHANGELOG.md} for the current behavior. The
$R$ values for TBR23 are from the April analysis and were not recomputed.}

\subsection{Mean cross-shore flow validation}
\label{sec:mean_flow_validation}

\statusnote{Historical method (May 2026). This is the test design behind the
mean-flow results in the Results section. It ran on 38 instruments at both 1-hour
and 17-minute segmentation. The 17-minute archive was deleted from the server on
2026-07-28, so the segmentation comparison cannot be rerun until a 17-minute L2 is
regenerated from the current L1 (\texttt{validation/archive/compare\_seglen\_phase2.m}).}

The segment-mean cross-shore velocity $\overline{u}$ was tested against three
hypotheses (wave-driven Stokes return flow, a constant per-instrument calibration
offset, and random Doppler noise) on every PUV deployment then in the catalog
(38 instruments at Torrey Pines, Los Pe\~nasquitos lagoon, Solana Beach, SIO Pier,
Catalina Island, and the Ruby2D cross-shore array).

\subsubsection{\texorpdfstring{Linear scaling against $H_s^2$}{Linear scaling against Hs squared}}
For each instrument the segment-mean cross-shore velocity was fit as
%
\begin{equation}
  \overline{u} = \alpha\, H_s^2 + \beta + \varepsilon,
  \label{eq:mean_flow_fit}
\end{equation}
%
restricted to segments with $H_s \geq \SI{0.2}{\meter}$. Stokes return-flow theory
(Longuet-Higgins, 1953) predicts a depth-averaged Eulerian compensation flow
$\overline{u} \approx -gH_s^2 / (16ch)$ for shallow water with $c = \sqrt{gh}$,
so $\alpha_\text{theory} = -g/(16ch)$, about $-0.018$ at $h = \SI{5}{\meter}$
and $-0.003$ at $h = \SI{15}{\meter}$ in \si{\meter\per\second} per
\si{\meter\squared} of $H_s^2$, with $\beta_\text{theory} = 0$. The hypotheses
predict different coefficients. Real return flow gives $\alpha$ matching theory
in sign and within a factor of order one, with $\beta$ centered on zero across
instruments. A constant offset gives $\alpha \approx 0$ with $\beta$ scattered at
the offset magnitude. Random noise leaves both coefficients unconstrained but
shrinking toward zero with longer averaging.

\subsubsection{Three robustness checks}
First, the same fit was run on the alongshore component $\overline{v}$, to test
whether $\alpha$ is specific to cross-shore physics or reflects storm-coherent
forcing (wind stress, atmospheric setup, alongshore gradients) that also
correlates with $H_s$. Second, the fit was restricted to
$H_s \geq \SI{1}{\meter}$ and the change in $R^2$ reported. Third, the alongshore
model was extended to
%
\begin{equation}
  \overline{v} = \alpha_{v0}\, H_s^2 + \alpha_{v1}\, \sin(2\theta_\text{rel})\, H_s^2 + \beta_v,
  \label{eq:wave_dir_check}
\end{equation}
%
where $\theta_\text{rel}$ is the wave incidence angle relative to shore normal,
taken from the L2 mean direction in the rotated frame. An error
$\theta_\text{err}$ in the shore-normal angle predicts
$\alpha_{v0} = \alpha_u \tan\theta_\text{err}$ with $\alpha_{v1} \approx 0$,
while a radiation-stress-driven longshore current predicts a dominant
$\alpha_{v1}$ with the sign of the local wave climate. The implied
$\theta_\text{err} = \arctan(\alpha_{v0}/\alpha_u)$ is an upper bound on any
rotation error.

\subsubsection{Auxiliary diagnostics}
(1)~Tidal modulation of $\overline{u}$ was measured by binning segment means
into 16 phase bins of band-passed depth (\SIrange{0.0019}{0.0093}{\hertz}, periods
of about \SI{30}{\minute} to \SI{3}{\hour}), stratified by $H_s$; a constant
offset has zero modulation in every stratum. (2)~Time series of $\overline{u}$
from independently calibrated instruments were correlated at common timestamps.
(3)~The Reynolds stress $\overline{u'w'}$, a segment mean of fluctuation
products rather than of the velocity itself, was correlated with
$\overline{u}$; independent noise in either path would push that correlation
toward zero.

\subsubsection{Segmentation independence}
The validation was run at 1-hour and 17-minute segmentation. Random noise
predicts that $\overline{u}$ estimates shrink toward zero as the averaging window
grows. Real wave-driven flow predicts $\alpha$ roughly independent of
segmentation, since it is a kinematic relation, with $|\beta|$ tightening
slightly toward zero.


\newpage
\section{Results}
\label{sec:results}
% >>> The current consolidated user guide is docs/PUV_Pipeline_Guide.tex (built to
%     PUV_Pipeline_Guide.pdf). This file is retained as source material / deep reference;
%     its L1 quality-control description predates the 2026-07 per-channel QC rework.
%     Start newcomers on the guide.
% PUV Pipeline: results section, \input by compile_sections.tex.
% HISTORICAL RECORD (status set 2026-10-15). Analyses run April-July 2026 on the
% catalog as it then stood; none recomputed. Retracted and superseded passages
% are marked with visible \statusnote boxes rather than deleted. Audit:
% docs/doc_audit_2026-10-15/audit_compile_sections.md.
\label{sec:results-status}

\statusnote{\textbf{Historical results, not recomputed.} Everything in this
section was computed between April and July 2026 on the catalog as it then stood:
33 to 38 instruments for most subsections, 17-minute segments for the earliest
(bulk agreement, spectral bias, Ruby2D, L3, TBR23 recovery), 61 records for the
July hypothesis tests. Since then the catalog has been rebuilt several times:
heading corrections (2026-05-14, 2026-07-27, 2026-10-07), \texttt{doffp},
coordinate and clock corrections (2026-07-27/28), the edge-gap fill fix and new
sediment physics (2026-10-09), and the full L2 to L4 rebuild for catalog
2026-10-14.1. Treat every number below as a record of what was found at the time.
Current versions of the model-comparison results are maintained in the
\texttt{Wave\_model\_validation\_paper} repository, where a reconciliation
against catalog 2026-10-14.1 was in progress on 2026-10-15.}

\subsection{PUV--MOP validation}

\subsubsection{Bulk wave parameter agreement}

Of 40 instruments across 20 deployments, 33 were processed successfully
(82.5\%), with 7 failures attributable to hardware issues (burial, bent
pipes, dead batteries). PUV-derived wave parameters showed strong
agreement with the MOP wave model across all four TBR23 instruments
(the summary table planned here was never drafted). For instruments unaffected by biofouling,
$H_s$ correlations exceeded $R^2 = 0.83$ with RMSE $<
\SI{0.06}{\meter}$ and biases of \SIrange{+0.002}{+0.008}{\meter}. One
instrument (MOP580 5m) exhibited lower agreement ($R^2 = 0.66$, bias $=
\SI{+0.037}{\meter}$), attributed to intermittent kelp fouling that
reduced the valid segment rate to 49\% compared with
\SIrange{85}{91}{\percent} for the remaining instruments. A small but
consistent positive $H_s$ bias was observed across all instruments, with
PUV-measured wave heights systematically exceeding shoaled MOP values.

% TODO: Insert Table \ref{tab:validation} with R2, RMSE, bias, skill, N
% for Hs, Tm, Ef across all instruments

\subsubsection{Frequency-dependent spectral bias}

The frequency-resolved spectral ratio
$S_\eta^\text{PUV}(f) / S_\eta^\text{MOP,shoaled}(f)$, averaged across
all matched hours, revealed that the positive $H_s$ bias was concentrated
at the swell spectral peak (\SIrange{0.04}{0.08}{\hertz}), where the
pressure transfer function $K_p > 0.9$ and the correction is minimal
(table planned, never drafted). The mid-swell band
(\SIrange{0.08}{0.14}{\hertz}) showed near-unity ratios ($0.99$--$1.03$),
and the high-swell band (\SIrange{0.14}{0.25}{\hertz}), where the
pressure correction is most aggressive ($K_p$ as low as $0.24$), also
showed only modest excess ($1.03$--$1.05$). The low-swell amplification
increased with decreasing depth, from a band-averaged ratio of $1.17$ at
\SI{8.3}{\meter} to $1.28$ at \SI{5.5}{\meter}. The overall excess
energy relative to the shoaled MOP spectrum was \SI{4.5}{\percent} at
\SI{8.3}{\meter} depth and \SI{8.7}{\percent} at \SI{5.5}{\meter}.

The concentration of excess energy at the spectral peak, where the
pressure correction factor is close to unity, indicates that the bias
does not originate from the pressure-to-surface-elevation transfer
function.

\statusnote{Superseded 2026-07. The apparent concentration of excess at the peak
came from comparing spectra on mismatched frequency grids (see the retraction
under ``MOP spectral peak broadening'' below). On the model's own frequency bins
the excess sits near the second harmonic, $f/f_p \approx 2$.}

% TODO: Insert Table \ref{tab:spectral_ratio} with band-averaged ratios
% and Kp ranges for each instrument

% Rewritten 2026-07-27. The earlier version of this section eliminated three
% mechanisms in favour of a fourth (MOP peak broadening) that was subsequently
% shown to be an artifact of comparing spectra on mismatched frequency grids.
% Re-tested on 72,948 hours / 61 records with
% validation/redo_hypothesis_elimination.m, the elimination inverted: nonlinear
% shoaling is the mechanism, bound long waves are its difference-frequency
% counterpart, directional narrowing is a real but separate bias, and peak
% broadening is eliminated. Full account:
% ../PUV_paper/docs/findings_hypothesis_elimination_2026-07-27.md
\subsubsection{Hypothesis testing}

\statusnote{Rewritten 2026-07-27; values partly superseded 2026-07-28. This
subsection replaced an April version that ruled out the first three mechanisms in
favor of the fourth. The values are from \texttt{validation/redo\_hypothesis\_elimination.m}
as run on 2026-07-27 (72,948 hours, 61 records). After the 2026-07-28 catalog
refresh, \texttt{Wave\_model\_validation\_paper/docs/findings\_refresh\_2026-07-28.md}
records these values as moved (not recomputed here): $\rho(\mathit{Ur}, \nu)$
$+0.332 \to +0.347$; $\nu$ in the most linear $H_s/h$ bin $1.000 \to 1.003$;
the ratio at $f/f_p = 2.1$ for $H_s/h > 0.12$, $1.338 \to 1.332$; bicoherence
ordering $-0.008 \to +0.017$ (lowest band) and $+0.285 \to +0.304$ (highest);
$\sigma_1$ ratio $1.156 \to 1.157$; matched-grid $Q_p$ ratio $1.008 \to 1.009$.
The verdicts did not change on that refresh. None of these has been recomputed
on catalog 2026-10-14.1.}

Four candidate mechanisms for the systematic positive bias were evaluated
(summary table planned, never drafted).

\paragraph{Nonlinear shoaling.}
Triad interactions during shoaling transfer energy from the spectral peak to its
super-harmonic near $2f_p$. The PUV/model energy ratio, formed on the model's
own frequency bins and binned by $f/f_p$, shows that transfer directly: a trough
of $0.94$--$0.95$ at $f/f_p = 1.3$--$1.5$ and a maximum of $1.34$ at
$f/f_p = 2.1$ for hours with $H_s/h > 0.12$, against $1.10$ for hours below it.
The discrepancy is organized by the Ursell number
($\rho = +0.332$, $N = \num{72948}$ hours across 61 records) rather than by
depth, which drops to $\rho = -0.114$ once $H_s/h$ is held fixed. Binned by
$H_s/h$, the spectral-width ratio $\nu^\text{PUV}/\nu^\text{MOP}$ rises
monotonically from $1.000$ below $H_s/h = 0.04$ to $1.100$ above $0.20$: the
model and the observations agree to three decimal places when the waves are
linear, and diverge as they cease to be. Bicoherence at the swell
self-interaction orders the harmonic excess \emph{within} fixed $H_s/h$ bands,
from $\rho = -0.008$ below $H_s/h = 0.06$ to $\rho = +0.285$ above $0.20$:
flat where triad coupling cannot operate, strengthening where it can.

An earlier version of this analysis correlated the Ursell number against the
\emph{low-swell} band ratio (\SIrange{0.04}{0.08}{\hertz}) and found
$R^2 = 0.004$. That test cannot detect the mechanism: nonlinear shoaling deposits
energy near $2f_p \approx \SI{0.18}{\hertz}$, not at the peak.

\paragraph{Directional narrowing.}
The first-harmonic directional spread $\sigma_1$, from the PUV cross-spectral
$a_1$ and $b_1$, is \emph{wider} than the MOP spread at swell frequencies, by a
median factor of $1.156$ across 61 records ($p = \num{2.6e-6}$). The sign is
opposite to the narrowing expected from refraction-driven focusing, and the
wider PUV spread is consistent with scattered energy at shallow depths that the
forward-propagating model does not carry. The spread ratio does not covary with
the discrepancy it was proposed to explain ($p = 0.30$ against the $H_s$ ratio,
$p = 0.54$ against $\nu$), so it is a separate bias in the model's directional
distribution rather than the mechanism behind the energy difference. Because
$\sigma_1$ depends on $r_1 = \sqrt{a_1^2 + b_1^2}$, it is invariant to
instrument rotation and unaffected by the heading corrections applied in 2026-07.

\paragraph{Bound long waves.}
Second-order difference-frequency forcing predicts bound infragravity variance
from the swell band via the Hasselmann (1962) / Sch\"affer--Madsen (1995) kernel.
The ratio of predicted-bound to observed-total IG variance rises with $H_s/h$
($\rho = +0.913$ across 61 record medians; $+0.758$ per hour) and crosses unity
near $H_s/h = 0.12$, above which second-order theory over-predicts the bound
component. Against the spectral-width discrepancy the same ratio gives
$\rho = +0.288$, falling to $+0.149$ once $H_s/h$ is controlled. Bound-wave
forcing and super-harmonic generation both arise from triad interactions among
the same swell components, so this collinearity is expected; the two are aspects
of one nonlinear response rather than competing explanations, and the data here
do not separate them.

An earlier version rejected this hypothesis partly on the grounds that the excess
energy sat at the swell peak rather than at group frequencies. That apparent
concentration at the peak was itself the grid artifact described below.

% =====================================================================
% RETRACTED 2026-07-25 — READ BEFORE REUSING ANYTHING IN THIS PARAGRAPH
% OR THE "Cross-deployment confirmation" SUBSECTION THAT FOLLOWS.
%
% The peak-broadening result below is an ARTIFACT of the comparison
% method, not a property of the MOP model. It was computed by
% interpolating MOP's ~20 coarse frequency bins UP onto the PUV's
% 3601-point grid (validation/analyze_spectral_shape.m:110) and then
% measuring half-power bandwidth and Qp on the fine grid. MOP bandwidths of
% 0.005-0.10 Hz against a multi-taper resolution bandwidth of 2*NW/T =
% 2.2e-3 Hz is a ~2-45x mismatch. (Stated as "74x" before 2026-07-27, which
% compared against the BIN SPACING of 2.78e-4 Hz; an NW = 4 estimate is
% oversampled 8x relative to what it resolves. The empirical 68.3% artifact
% is unaffected -- it was measured, not derived from the ratio.) Linear interpolation between coarse bin centres
% manufactures a broad smooth peak and depresses int(S^2 df), the
% numerator of Qp. Both metrics move in the direction of the claim.
%
% Measured against a known-zero truth: a PUV spectrum compared against a
% degraded copy of ITSELF through that same path yields 68.3% apparent
% bandwidth narrowing and a Qp ratio of 1.225 (TBR23/MOP586_7m, 1382
% segments). That EXCEEDS the entire effect reported here. The m0 ratio
% of 0.9954 confirms it is a shape artifact, not an energy error.
%
% On a resolution-matched grid (both products reduced to the model's
% native bins first) the effect is absent and changes sign across sites:
% Qp ratio 1.093 Torrey, 1.041 Cardiff, 0.930 Coronado, 0.999 Torrey
% offshore. The matched metric recovers injected broadening at 95-96%
% efficiency down to r = 1.02, so this is an informative null.
%
% The factor-of-3 disagreement between the 16.7% here and the 36-78% in
% cross_deployment_shape_summary.mat is explained: the apparent effect
% scales with the PUV estimator's resolution (Welch/17-min vs
% multi-taper/1-hour), not with anything physical.
%
% Replacement code:  validation/compare_shape_matched.m
%                    validation/run_matched_shape_sweep.m
% Regression tests:  validation/test_resolution_artifact.m
%                    validation/test_shape_metric_sensitivity.m
%                    validation/test_shoal_bin_synthetic.m
% Full write-up:     PUV_paper/docs/findings_resolution_artifact_2026-07-24.md
%
% The hypothesis-elimination prose above was REWRITTEN on 2026-07-27 to reflect
% the redone analysis. This banner is retained because the paragraph it guards
% (the peak-broadening result) is still present as the record of what was
% originally claimed. See
% validation/redo_hypothesis_elimination.m and
% ../PUV_paper/docs/findings_hypothesis_elimination_2026-07-27.md:
%
%   H1 nonlinear shoaling    was "ruled out" -> IS THE MECHANISM
%       rho(Ursell,nu) = +0.332; nu = 1.000 in the most linear Hs/h bin;
%       excess localised at f/fp = 2.1 (1.338 at high Hs/h); bicoherence
%       orders it within Hs/h bands (+0.285 above Hs/h = 0.20)
%   H2 directional narrowing was "ruled out" -> REAL but SEPARATE bias
%       sigma_1 ratio 1.156, p = 2.6e-6, yet uncorrelated with the
%       discrepancy (p = 0.30, 0.54). Report it; it is not the cause.
%   H3 bound long waves      was "ruled out" -> REAL, collinear with H1
%       +0.913 per record vs Hs/h, but only +0.149 against nu after
%       controlling for Hs/h
%   H4 peak broadening       was THE WINNER  -> ELIMINATED (artifact)
%
% The three sections below that assert these were "ruled out" must be
% REWRITTEN, not annotated. The original analysis did not merely get the
% answer wrong -- it eliminated the correct one.
% =====================================================================

\statusnote{\textbf{Retracted 2026-07-24.} The result in the next two paragraphs
is an artifact of the comparison method, not a property of the MOP model. It was
computed after interpolating the model's roughly 20 frequency bins onto the PUV
grid (\texttt{validation/archive/analyze\_spectral\_shape.m}), which manufactures a
broad, smooth peak and lowers $\int S^2\,df$. Comparing a PUV spectrum with a
degraded copy of itself through the same path gives a $Q_p$ ratio of 1.225 and
68\% apparent narrowing (TBR23/MOP586\_7m, 1382 segments), more than the whole
effect reported here. On a resolution-matched grid the effect is absent; see
``Cross-deployment confirmation'' below. The paragraphs are kept as the record of
what was claimed. Quote nothing from them.}

\paragraph{MOP spectral peak broadening.}

The Goda peakedness parameter $Q_p$ was systematically higher for PUV
spectra than for shoaled MOP spectra (median $Q_p = 1.88$ vs.\ $1.75$,
ratio $1.07$). The PUV half-power spectral bandwidth was
\SI{16.7}{\percent} narrower than MOP (\SI{0.039}{\hertz} vs.\
\SI{0.047}{\hertz}), and the PUV peak spectral density was
\SI{14.4}{\percent} higher. The peakedness difference
$\Delta Q_p = Q_p^\text{PUV} - Q_p^\text{MOP}$ showed the strongest
correlation with $H_s$ amplification of any tested mechanism
($R = 0.44$, $R^2 = 0.20$). Hours when the PUV spectrum was more peaked
relative to MOP corresponded to hours with larger positive $H_s$ bias.

These results indicate that the MOP spectral wave model introduces
numerical broadening of the spectral peak during propagation from the
initialization depth (\SI{10}{\meter}) to the instrument depth. The
in-situ PUV measurement captures the true, narrower spectral peak shape,
resulting in higher peak energy density and thus higher $H_s$ for a
given total variance. This spectral broadening is a recognized property
of discretized spectral wave models (cited in the April draft by key only: \texttt{rogers2002}, \texttt{booij1999})
and has not been previously quantified using co-located PUV observations
in this depth range.

\subsubsection{Cross-deployment confirmation}

\statusnote{Rewritten after the 2026-07-24 retraction. The April version of this
subsection reported the broadening pattern as reproduced across 11 instruments.}

% RETRACTED 2026-07-25 along with the paragraph above -- see that banner.
% Two separate problems with this subsection:
%   (1) The pattern it reports is the resolution artifact, which is present
%       in every instrument by construction. Reproducibility across
%       deployments is exactly what an artifact produces; consistency of
%       sign was read as evidence of robustness when it was evidence of a
%       shared upstream computation.
%   (2) "11 instruments" is stale even on its own terms. The saved result
%       (outputs/validation/cross_deployment_shape_summary.mat) holds 33
%       instrument-records, and the current catalog is 65. The 11 figure
%       also appears in PIPELINE_NOTES.md and docs/puv_pipeline_slides.tex.

On a resolution-matched grid the peak-shape difference is absent across the
catalog: the median $Q_p$ ratio is $1.008$ over 62 instrument-records
($p = 0.44$), with an interquartile range of $0.980$--$1.039$ that spans unity.
The legacy interpolated-grid calculation returns $1.203$ on the same records,
within round-off of the $1.225$ obtained by comparing a PUV spectrum against a
degraded copy of itself.

Consistency of the original pattern across instruments was read as evidence that
it was physical. It is not: the artifact is a property of the comparison, so it
appears in every record by construction, and agreement across deployments at
matched depths is what a shared upstream error produces.

% TODO: Insert Table \ref{tab:hypotheses} summarizing hypothesis,
% test metric, result, and verdict for each of the four mechanisms

\subsubsection{Pipeline cross-validation: Ruby2D head-to-head}
\label{sec:ruby2d-validation}

\statusnote{April 2026, 17-minute segments. The new-pipeline side of this
comparison used 2048-sample segments ($\Delta f = \SI{9.77e-4}{\hertz}$) and the
L2 code of that date. The later, isolated test of each legacy fault
(\texttt{docs/legacy\_pipeline\_defects.pdf}, 3,713 segments, 10 records) finds
the band-averaged spread error reaching \SI{17}{\degree} at the 95th percentile
and the $b_2$ moment, which radiation stress is built from, unreliable. The
\SI{1}{\degree} median spread offset reported here is consistent with that, but
the conclusion that prior Ruby2D statistics need no revision applies to $H_s$,
$T_p$ and mean direction only. The legacy source as it now stands sets \texttt{nfft} = 2400 for both auto-
and cross-spectra, with comments recording a change from 7200. The archived
Ruby2D L2 reports $T_p$ at 1200, 1800 and \SI{3600}{\second}, the first three
bins of a \SI{2.78e-4}{\hertz} grid, so it was built before that edit.}

To verify that the multi-taper L2 implementation does not introduce
biases relative to an established processing pipeline, we reprocessed
the Ruby2D MOP582\_6m record (Torrey Pines, October 2021 to February 2022,
$\sim$3,100 hours) and compared it segment-by-segment against the
archived L2 product produced by an independent prior pipeline
(\texttt{beachpeeps/PUV\_Processing}). The prior
pipeline uses a single-periodogram estimator (Hann window, full
60-minute segments, $\Delta f = \SI{2.78e-4}{\hertz}$) and applies a
fixed cap of $K_p^2 \le 10$ to the pressure transfer function. The legacy
record has 3,095 60-min segments; the new pipeline produces 10,903 17-min
segments over the same record. Segments were matched on midpoint
nearest neighbor with a tolerance of 30~minutes, yielding 2,322 matched
pairs (75\%). The remaining 25\% are hours that the new pipeline rejects
under the segment-validity criterion (\texttt{nanMaxFrac} = 0.10) but
that the prior pipeline retains by zero-padding small NaN gaps.

\begin{table}[htbp]
  \centering
  \caption{Ruby2D MOP582\_6m: bulk wave parameter agreement between
    the new multi-taper pipeline and the legacy archived L2. $T_p$ statistics
    are computed after restricting both estimates to the SS band
    (\SIrange{4}{25}{\second}); without this filter, three segments on
    January 15, 2022 in the prior pipeline report $T_p$ at the IG/DC
    bin ($T_p = 1200$, 1800, 3600~\si{\second}) because the prior
    pipeline's peak picker is not band-limited.}
  \label{tab:ruby2d}
  \begin{tabular}{@{}lccccc@{}}
    \toprule
    Metric & Legacy (med.) & Multi-taper (med.) & Bias (M$-$L) & RMS & $R^2$ \\
    \midrule
    $H_s$ (\si{\meter})        & 0.771  & 0.757  & $-0.009$ & 0.053 & 0.981 \\
    $T_p$ (\si{\second})       & 13.48  & 13.30  & $-0.02$  & 1.48  & 0.761 \\
    Direction SS (\si{\degree})    & $-7.3$ & $-7.2$ & $+0.2$   & 1.2   & 0.927 \\
    Spread SS (\si{\degree})       & 14.7   & 15.8   & $+1.0$   & 1.7   & 0.877 \\
    \bottomrule
  \end{tabular}
\end{table}

Bulk wave parameter agreement is excellent across all four metrics
(Table~\ref{tab:ruby2d}; Figure~\ref{fig:ruby2d_scatter}). $H_s$ agrees
to 5~cm RMS with a $-9$~mm median bias
($R^2 = 0.98$), and direction agrees to within
1.2$^\circ$~RMS ($R^2 = 0.93$). The directional spread is
1$^\circ$ wider in the new pipeline on average; this is the largest
systematic offset observed and remains within the segment-to-segment
scatter. The new $T_p$ shows the discrete quantization bands expected from
$\Delta f = 0.000976$~\si{\hertz} versus the prior pipeline's
$0.000278$~\si{\hertz}, but no systematic bias.

\begin{figure}[htbp]
  \centering
  \includegraphics[width=0.85\textwidth]{Ruby2D/ruby2d_582_6m_scatter.png}
  \caption{Bulk parameter agreement between the new multi-taper pipeline
    (vertical axis) and the legacy archived L2 (horizontal axis) for
    Ruby2D MOP582\_6m, 2,322 matched 60-minute segments,
    October 2021 to February 2022. The dashed line is 1:1.}
  \label{fig:ruby2d_scatter}
\end{figure}

Per-segment surface elevation spectra agree closely in shape
(Figure~\ref{fig:ruby2d_spectra}). At a representative median segment
($H_s = 0.77$~\si{\meter}, November 27, 2021) and at a storm peak
($H_s = 2.78$~\si{\meter}, December 14, 2021, bimodal sea state), the
integrated $H_s$ values from the two pipelines agree to within 4\%.
The visible difference between the spectra is the chi-square noise
characteristic of a single-periodogram estimator in the prior pipeline,
versus the variance reduction provided by averaging $K = 7$ DPSS tapers
in the new pipeline. The underlying peak shapes and locations agree.

\begin{figure}[htbp]
  \centering
  \includegraphics[width=0.95\textwidth]{Ruby2D/ruby2d_582_6m_spectra.png}
  \caption{Representative surface elevation spectra for Ruby2D
    MOP582\_6m. Top: median segment ($H_s = 0.77$~\si{\meter}). Bottom:
    storm peak ($H_s = 2.78$~\si{\meter}). Left column: log scale,
    full SS band. Right column: linear scale, swell band detail. The
    multi-taper estimates (blue) are smoother than the prior single-periodogram
    estimates (red); peak locations and total variance agree.}
  \label{fig:ruby2d_spectra}
\end{figure}

This cross-validation establishes two things. First, the multi-taper
implementation reproduces an independent prior pipeline to within
instrument noise on bulk wave parameters, confirming that the
estimator change is a refinement rather than a methodological break.
Prior published Ruby2D bulk wave statistics do not require revision.
Second, the consistency-of-windows fix in the new cross-spectral
calculations (auto- and cross-spectra both use the same DPSS tapers,
unlike the prior pipeline which mixed Hanning and Hamming) does not
produce the order-20\% directional errors that the bug analysis
suggested were possible at this site: the median impact at
Ruby2D MOP582\_6m is approximately $1^\circ$ in spread.

\subsubsection{Site-specific amplification at SIO Pier}

The SIO Pier instruments (MOP 511, ~7m depth) exhibited dramatically
larger PUV-to-MOP spectral amplification than all other sites. The
peak spectral density ratio ranged from 2.0 to 2.7 across four
deployment periods (SIO24A through SIO25E), with corresponding $H_s$
amplification of 29--36\%, far exceeding the 1--11\% observed at
Torrey Pines, Solana Beach, and Los Pe\~nasquitos sites at comparable
depths.

Investigation confirmed that the MOP frequency grid is identical across
all stations (20 bins, 0.04--0.40~Hz), ruling out a resolution artifact.
The MOP model total energy at 10~m depth is comparable between SIO
($H_s = 0.62$~m) and Torrey Pines ($H_s = 0.61$~m), yet the in-situ
PUV at SIO measures approximately 2.5$\times$ the peak spectral density
predicted by the shoaled MOP spectrum.

The SIO Pier PUV is located near the head of Scripps Submarine Canyon,
where canyon-edge refraction is known to create wave energy focusing
(the April draft cites these by bibliography key only, \texttt{munk1948} and
\texttt{shepard1950}; the references were not identified or checked in the 2026-10
revision). The canyon bifurcates nearshore,
and wave focusing at the canyon head has been documented at Black's
Beach on the adjacent fork. The MOP model, which propagates spectra
along a relatively coarse alongshore grid, does not resolve this
fine-scale bathymetric focusing. The PUV measurements capture the true,
focused wave field, resulting in the observed amplification.

\statusnote{The comparison with ``universal spectral peak broadening'' in the next
paragraph refers to the result retracted on 2026-07-24. The SIO amplification
itself was measured on the then-current spectra; whether it survives the
resolution-matched comparison was not checked in this document.}

This site-specific effect is distinct from the universal spectral peak
broadening described above. While broadening produces a modest
($\sim$5--10\%) Hs bias that is consistent across all sites, the
canyon focusing at SIO produces an order-of-magnitude larger
amplification that is unique to that location.

\subsection{L3 wave forcing characterization}

\statusnote{April 2026, 33 instruments, 17-minute segments. The Rouse numbers
below were computed from the oscillatory wave-stress amplitude, which overstated
suspension; since 2026-08-13 the Rouse number uses the mean combined
wave--current stress, and the corrected values are larger (median 41.5 against
4.8 at TBR23 MOP580\_5m). The bedload conclusion is unchanged in sign, but the
quoted range is obsolete. Shields values used $D_{50} = \SI{0.25}{\milli\meter}$
everywhere; TBR23 now uses measured grain sizes. The ``UTC clock convention''
was confirmed for the instruments tested here; the pre-2019 file-name clocks
were later found to be Pacific local time.}

\subsubsection{Energy band decomposition}

Frequency-band analysis revealed clear seasonal patterns across all
deployments. Winter deployments (TOR23W, SOL23, TOR24W, SOL24) were
swell-dominated (65--90\% of total energy flux in the 0.04--0.10~Hz
band), while summer deployments (TBR23, SIO25B--D) showed
approximately equal swell and sea contributions. The energy flux
consistency check confirmed exact agreement between the sum of band
fluxes and the total L2 energy flux ($F_\mathrm{total}/F_\mathrm{L2}
= 1.0000$) for all 33 instruments.

\subsubsection{Storm characterization with MOP gap-filling}

Storm detection identified 0 events during the calm summer TBR23
deployment and up to 11 events during winter TOR23W, including the
major December 28, 2023 storm sequence ($H_s$ up to 3.6~m at 7~m
depth, sustained over 74~hours). MOP gap-filling proved essential:
the SOL23 MOP651 7m instrument detected 0 storms from PUV data
(55\% valid segments) but MOP identified 3 events during data gaps,
including an 86-hour storm with $H_s = 3.09$~m. Across all
deployments, MOP supplemented 1--8 additional storm events per
instrument that were missed due to PUV data gaps during the most
energetic conditions.

\subsubsection{Transport proxy patterns}

The Shields parameter showed clear depth and seasonal dependence.
Median Shields numbers decreased from 0.24 at 5.6~m to 0.065 at
15.5~m depth (TOR23W MOP586 cross-shore array), with corresponding
mobilization fractions decreasing from 100\% to 78\%. Winter
deployments showed approximately 3$\times$ the cumulative bottom
energy flux of summer deployments at equivalent depths. The Rouse
number indicated bedload-dominated transport (median $P = 2.8$--6.0)
across all sites and seasons, with the lowest values (most
suspension-capable) at the shallowest winter sites.

\subsubsection{Tidal and subtidal currents}

Tidal currents were small (typically $<$0.05~m/s) relative to
wave-driven subtidal flows. The subtidal cross-shore residual was
consistently negative (offshore-directed) across all sites and
seasons, with magnitudes of 0.004--0.023~m/s, consistent with
wave-driven undertow. Tidal depth predictions from the Scripps Pier
NOAA gauge matched PUV depth anomalies at $R > 0.99$ for instruments
with continuous records, confirming the UTC clock convention and
validating the pressure-to-depth conversion.

\subsection{TBR23 post-storm recovery dynamics}

\statusnote{April 2026. This analysis moved to the Paper\_1 repository, which
holds the current numbers. Note for the MOP580\_5m exclusion below: on
2026-05-14 that record's L1 heading was found to be \SI{180}{\degree} wrong (a
field-sheet entry of \SI{270.74}{\degree} for \SI{90.74}{\degree}) and was
corrected. A \SI{180}{\degree} heading error reverses the sign of $u$, which by
itself produces negative velocity skewness and anti-correlation of
$\overline{u}$ with neighboring instruments, the two symptoms attributed here to
kelp fouling. The fouling attribution of those symptoms is therefore doubtful.
The low valid-segment rate (49\%) does not depend on heading and was not
re-examined.}

\subsubsection{Observed morphological change}

Jetski bathymetric surveys (22 surveys during the deployment period)
revealed 25--39~cm of net erosion at the PUV instrument locations
over the May--August 2023 deployment, with the 5~m sites eroding
more than the 7~m sites. Pre-deployment context from the MOPS survey
archive showed that the winter 2023 storm sequence had deposited
approximately 85~cm of sand at the 5~m depth contour relative to the
prior year, placing the bed well above its long-term equilibrium
profile. The TBR23 deployment thus captured the post-storm
readjustment phase rather than active storm erosion.

\subsubsection{Competing transport mechanisms}

Analysis of the PUV velocity moments revealed two competing
cross-shore transport mechanisms operating simultaneously during the
recovery period:

\paragraph{Onshore: velocity skewness.}
Velocity skewness $S_k$ was positive at all three reliable instruments
(median $S_k = +0.05$ to $+0.12$), indicating peaked crests and flat
troughs characteristic of shoaling waves. Skewness increased with
decreasing water depth, following the Ruessink et al.\ (2012)
parameterization based on the Ursell number. The skewness-driven
bedload flux $\langle u |u|^2 \rangle$ was directed onshore,
representing the wave-driven recovery mechanism.

\paragraph{Offshore: undertow.}
The subtidal mean cross-shore current was consistently offshore-directed
($\overline{u} = -0.006$ to $-0.012$~m/s) at all three instruments,
representing wave-driven return flow (undertow). Undertow magnitude
increased with wave height and was modulated by tidal phase, with
stronger offshore flow at high tide when more wave energy reached the
instrument depths.

\subsubsection{Depth-dependent transport balance}

The ratio of onshore skewness flux to offshore undertow flux showed a
clear depth dependence exploiting the tidal modulation of water depth
at each fixed instrument location ($\sim$2.5~m tidal range). At the
shallowest site (MOP586 5m, median depth 5.5~m), the skewness flux
exceeded the undertow flux by a factor of 3--4, indicating net onshore
transport and active recovery. At the 7~m sites (median depth 8.3~m),
the undertow flux slightly exceeded the skewness flux, indicating near-
equilibrium or weak offshore transport.

The crossover depth where onshore and offshore fluxes balanced
was approximately 6--7~m, defining a \emph{recovery boundary} below
which the post-storm deposit was actively being redistributed
shoreward. This is consistent with the survey observations: the
MOP586 5m site exhibited a 25~cm recovery pulse in mid-July 2023 that
was not observed at the deeper instruments.

One instrument (MOP580 5m) was excluded from the transport analysis
due to intermittent kelp fouling that produced anomalous negative
velocity skewness (92\% of segments) inconsistent with the other
instruments and with wave theory.

\subsection{Mean cross-shore flow validation}
\label{sec:mean_flow_results}

\statusnote{May 2026 (run of 2026-05-05). This run predates the 2026-05-14
heading corrections: TOR24S/MOP586\_7m entered it with its velocity sign
reversed, and TBR23/MOP580\_5m was excluded as kelp-fouled for symptoms a
reversed heading also produces (see the note on the TBR23 recovery section). It
also predates the 2026-07 and 2026-10 heading changes and the 2026-10 edge-gap
fill fix. The fractions of instruments with the predicted sign of $\alpha$ are
the numbers most exposed. Not rerun since.}

The validation framework defined in
Section~\ref{sec:mean_flow_validation} was applied across the
2023+ catalog (33 instruments at four sites: 19 at Torrey Pines,
4 at Los Pe\~nasquitos lagoon, 5 at Solana Beach, 5 at SIO Pier).
The same fits were repeated at the legacy 17-minute segmentation
to provide a segmentation-independence check; pre-2023 deployments
(Catalina Island, Ruby2D Torrey cross-shore array; total 5
additional instruments) were added at 1-hour only. The analysis
excludes the kelp-fouled TBR23 MOP580 5m record and the
SIO25C partial deployment ($N = 61$ segments, $R^2 = 0.001$).

\subsubsection{Cross-deployment headline}
The results decisively favor the real-Stokes-return-flow hypothesis
(headline table planned, never drafted; values are in
\texttt{docs/mean\_flow\_validation\_plan.md}, which is kept out of the public
repository). The intercept distribution
is centered on zero across both segmentations (median $|\beta| =
\SI{0.6}{\milli\meter\per\second}$ at 17-min,
$\SI{2.8}{\milli\meter\per\second}$ at 1-hour), and 100\,\% of
17-min and 94\,\% of 1-hour intercepts fall inside a $\pm$2~cm/s
band that envelopes the manufacturer-quoted single-sample velocity
accuracy at the deployed range setting. The slope distribution
clusters near the Stokes prediction
($\alpha/\alpha_\text{theory} = +0.53$ at 17-min, $+0.71$ at
1-hour), with 66--81\,\% of instruments showing the predicted sign
of $\alpha < 0$ (offshore-directed undertow). Going from 17-min to
1-hour averaging \emph{strengthens} the slope-vs-theory match and
the correct-sign fraction; random instrument noise would do the
opposite.

% TODO: Insert Table \ref{tab:mean_flow_headline}
%   columns: metric | 17-min | 1-hour
%   rows: median |beta|; |beta|<2cm/s fraction; median alpha/alpha_th;
%         alpha<0 fraction; median R^2; median tidal modulation amplitude

\subsubsection{Site decomposition}
The Torrey Pines site (19 instruments across 5 deployments and
depths from 4.8 to 15.5~m) gives the cleanest theory match: 83\,\%
have the correct $\alpha$ sign and median $\alpha/\alpha_\text{theory}
= +0.77$ at 17-min. The Los Pe\~nasquitos lagoon site agrees
($+0.88$ ratio, 75\,\% correct sign at 17-min). At Solana Beach
and SIO Pier, the 17-min slopes are mixed in sign with low $R^2$,
but both sites \emph{flip} to predominantly negative $\alpha$ when
re-fit at 1-hour (Solana: 20\,\% $\to$ 100\,\% correct sign;
SIO Pier: 40\,\% $\to$ 80\,\%). This is interpreted as longer
averaging exposing a wave-driven signal that the 17-min sampling
variance was partially obscuring; the noise hypothesis predicts
the opposite shift.

\subsubsection{Tidal-phase modulation}
Across all instruments, the median peak-to-peak modulation
amplitude of $\overline{u}$ at $H_s \geq \SI{1}{\meter}$ is
\SI{2.6}{\centi\meter\per\second} (1-hour) and
\SI{2.8}{\centi\meter\per\second} (17-min), with peaks reaching
6--\SI{8}{\centi\meter\per\second} at multiple Torrey instruments.
A constant instrumental offset has zero modulation amplitude in
any $H_s$ stratum by definition; the observed amplitudes are 30--80
times the segment-mean random-noise floor and several times larger
than the largest fitted intercepts.

\subsubsection{Cross-instrument and Reynolds-stress consistency}
Time correlations of $\overline{u}(t)$ between independently-
calibrated Vectors at common timestamps were positive across all
non-fouled instrument pairs in TBR23, reaching $R = +0.57$ between
the two 7-m records. The kelp-fouled MOP580 5m record showed
\emph{anti}-correlation with all peers ($R = -0.33$ to $-0.48$),
providing an internal consistency check against the
fouling-flag exclusion. The Reynolds stress component
$\overline{u'w'}$, derived from the same record but via an
independent processing path, correlated positively with
$\overline{u}$ at all four TBR23 instruments
($R = +0.24$ to $+0.63$).

\subsubsection{Wave-direction discrimination}
The extended alongshore model
(Eq.~\eqref{eq:wave_dir_check}) was decisive against the
shore-normal-rotation-error hypothesis. Across 30 instruments at
1-hour segmentation, the wave-direction-dependent coefficient
$\alpha_{v1}$ dominated the wave-direction-independent coefficient
$\alpha_{v0}$ in magnitude in 90\,\% of records, with the 95\,\%
confidence interval on $\alpha_{v1}$ excluding zero in 83\,\% of
records. The sign of $\alpha_{v1}$ was positive in 80\,\% of cases,
corresponding to NW swells driving southward longshore current and
S swells driving northward longshore current, the classical
radiation-stress signature for the WNW-facing San Diego County
coast. Reading $\alpha_{v0}/\alpha_u$ as an upper-bound implied
shore-normal angle error places that bound at a few degrees,
consistent with the precision of the CDIP MOP shore-normal
metadata.

\subsubsection{Verdict}
The cross-deployment dataset is incompatible with both the
constant-offset and the random-noise hypotheses, and matches the
Stokes return-flow prediction in sign, magnitude, and segmentation-
length response. Random Doppler noise reduces the segment-mean
noise floor to $\sim$0.4~mm/s at 17-min and $\sim$0.24~mm/s at
1-hour averaging, well below any of the slopes, modulation
amplitudes, or cross-instrument shared variance reported above. The
small (mm/s-scale) wave-direction-independent residuals visible in
the alongshore fit are consistent with site-specific Eulerian
residual currents and any small remaining shore-normal-rotation
error; they are not the dominant component of either velocity
component's response to wave forcing.


\end{document}
